% vocabulary of PDDL
\newcommand{\pddlTask}{\ensuremath{\Pi}\xspace}
\newcommand{\pddlDomain}{\ensuremath{\mathcal{D}}\xspace}
\newcommand{\pddlInstance}{\ensuremath{\mathcal{P}}\xspace}
\newcommand{\pddlObjects}{\ensuremath{O}\xspace}
\newcommand{\pddlInitial}{\ensuremath{I}\xspace}
\newcommand{\pddlGoal}{\ensuremath{G}\xspace}
\newcommand{\pddlState}{\ensuremath{S}\xspace}

\newcommand{\pddlPredicates}{\ensuremath{\mathsf{N}_\mathsf{P}}\xspace}
\newcommand{\pddlActions}{\ensuremath{A}\xspace}
\newcommand{\pddlDerRules}{\ensuremath{\text{Der}}\xspace}

\newcommand{\pddlPredicatesC}{\ensuremath{P_C}\xspace}
\newcommand{\pddlPredicatesR}{\ensuremath{P_R}\xspace}



\newcommand{\pddlActParam}{\ensuremath{\vec{x}}\xspace}
\DeclareMathOperator{\pddlActName}{name}
\DeclareMathOperator{\pddlActPre}{pre}
\DeclareMathOperator{\pddlActAdd}{add}
\DeclareMathOperator{\pddlActDel}{del}
\DeclareMathOperator{\pddlActionReq}{actionReqs}
\DeclareMathOperator{\pddlQReq}{qreq}

% vocabulary of ontology
\newcommand{\Ont}{\ensuremath{\mathcal{O}}\xspace}
\newcommand{\SwrlRule}{\ensuremath{s}\xspace}

\newcommand{\N}{\ensuremath{\mathsf{N}}\xspace}
\newcommand{\NC}{\ensuremath{\mathsf{N}_\mathsf{C}}\xspace}
\newcommand{\NR}{\ensuremath{\mathsf{N_R}}\xspace}
\newcommand{\NI}{\ensuremath{\mathsf{N_I}}\xspace}

\newcommand{\NRobject}{\ensuremath{\mathsf{N_R^{object}}}\xspace}
\newcommand{\NRdata}{\ensuremath{\mathsf{N_R^{data}}}\xspace}

\newcommand{\NIC}{\ensuremath{\mathsf{N_I^{C}}}\xspace}
\newcommand{\NIR}{\ensuremath{\mathsf{N_I^{R}}}\xspace}

\newcommand{\ALCI}{\ensuremath{\mathcal{ALCI}}\xspace}

\DeclareMathOperator{\antecedent}{antecedent}
\DeclareMathOperator{\consequent}{consequent}


% vocabulary of TOMASys
\DeclareMathOperator{\FunctionTomasys}{Function}
\DeclareMathOperator{\FunctionsTomasys}{Functions}
\DeclareMathOperator{\FunctionDesign}{FunctionDesign}
\DeclareMathOperator{\FunctionDesigns}{FunctionDesigns}
\DeclareMathOperator{\Component}{Component}
\DeclareMathOperator{\QAtype}{QAtype}
\DeclareMathOperator{\QAtypes}{QAtypes}
\DeclareMathOperator{\QAvalue}{QAvalue}
\DeclareMathOperator{\Objective}{Objective}
\DeclareMathOperator{\Objectives}{Objectives}
\DeclareMathOperator{\FunctionGrounding}{FunctionGrounding}
\DeclareMathOperator{\FunctionGroundings}{FunctionGroundings}
\DeclareMathOperator{\solvesF}{solvesF}
\DeclareMathOperator{\requiredBy}{requiredBy}
\DeclareMathOperator{\isQAtype}{isQAtype}
\DeclareMathOperator{\hasQAestimation}{hasQAestimation}
\DeclareMathOperator{\hasNFR}{hasNFR}
\DeclareMathOperator{\typeF}{typeF}
\DeclareMathOperator{\solvesO}{solvesO}
\DeclareMathOperator{\typeFD}{typeFD}
\DeclareMathOperator{\hasQAvalue}{hasQAvalue}
\DeclareMathOperator{\requiresO}{requiresO}
\DeclareMathOperator{\cStatus}{c\_status}

\DeclareMathOperator{\fdRealisability}{fd\_realisability}

\DeclareMathOperator{\hasValue}{hasValue}
\DeclareMathOperator{\qaHasValue}{qa\_has\_value}
\DeclareMathOperator{\qaComparisonOperator}{qa\_comparison\_operator}
\DeclareMathOperator{\qaCritical}{qa\_critical}
\DeclareMathOperator{\fdEfficacy}{fd\_efficacy}
\DeclareMathOperator{\fgStatus}{fg\_status}
\DeclareMathOperator{\oImprove}{o\_always\_improve}
\DeclareMathOperator{\oStatus}{o\_status}
\DeclareMathOperator{\oUpdatable}{o\_updatable}


%vocabulary new tomasys
\DeclareMathOperator{\Action}{Action}
\DeclareMathOperator{\fdBetterUtility}{fdBetterUtility}
\DeclareMathOperator{\functionGrounding}{functionGrounding}
\DeclareMathOperator{\requiresC}{requiresC}
\DeclareMathOperator{\requiresF}{requiresF}
\DeclareMathOperator{\fActive}{f\_active}
\DeclareMathOperator{\fdUtility}{fd\_utility}

% instances example
\DeclareMathOperator{\fdSpiralMed}{fd\_spiral\_medium}
\DeclareMathOperator{\fdSpiralLow}{fd\_spiral\_low}
\DeclareMathOperator{\fdSpiralHigh}{fd\_spiral\_high}
\DeclareMathOperator{\aSearchPipelineOwl}{a\_search\_pipeline}
\DeclareMathOperator{\fGenerateSearchPath}{f\_generate\_search\_path}
\DeclareMathOperator{\fMaintainMotion}{f\_maintain\_motion}
\DeclareMathOperator{\fFollowPipeline}{f\_follow\_pipeline}

\DeclareMathOperator{\Domain}{Domain}
\DeclareMathOperator{\Range}{Range}

\DeclareMathOperator{\varC}{?c}
\DeclareMathOperator{\varFD}{?fd}
\DeclareMathOperator{\varFDtwo}{?fd2}
\DeclareMathOperator{\varFDbetter}{?fd\_better}
\DeclareMathOperator{\varF}{?f}
\DeclareMathOperator{\varU}{?u}
\DeclareMathOperator{\varUtwo}{?u2}
\DeclareMathOperator{\varUbetter}{?u\_better}
\DeclareMathOperator{\varQA}{?qa}
\DeclareMathOperator{\varQAvalue}{?qav}

\DeclareMathOperator{\fdUnground}{fd\_unground}

\DeclareMathOperator{\instanceFD}{fd}
\DeclareMathOperator{\instanceF}{f}
\DeclareMathOperator{\instanceA}{a}

\DeclareMathOperator{\instanceTrue}{true}
\DeclareMathOperator{\instanceFalse}{false}

% vocabulary for TACA transformation
\DeclareMathOperator{\bestFD}{bestFD}
\DeclareMathOperator{\validAndBest}{valid}
\DeclareMathOperator{\functionalReq}{functionalReq}

% PlanSys algorithms
\DeclareMathOperator{\algE}{E}
\DeclareMathOperator{\algEff}{Eff}
\DeclareMathOperator{\algDerive}{Derive}
\DeclareMathOperator{\algGraph}{Graph}
\DeclareMathOperator{\algRoots}{Roots}
\DeclareMathOperator{\algS}{S}
\DeclareMathOperator{\algC}{C}
\DeclareMathOperator{\algP}{P}
\DeclareMathOperator{\algNodes}{Nodes}
\DeclareMathOperator{\algSCC}{SCC}
\DeclareMathOperator{\algSCCs}{SCCs}




% some useful comments
\newcommand{\hide}[1]{}



\newcommand{\gustavo}[1]{\todo[inline,color=cyan!30]{{\footnotesize Gustavo: #1}}}
\newcommand{\tobias}[1]{\todo[inline,color=orange!30]{{\footnotesize Tobias: #1}}}
\newcommand{\ebj}[1]{\todo[inline,color=red!30]{{\footnotesize Einar: #1}}}
\newcommand{\carlos}[1]{\todo[inline,color=green!30]{{\footnotesize Carlos: #1}}}
\newcommand{\francisco}[1]{\todo[inline,color=magenta!30]{{\footnotesize Francisco: #1}}}

\newcommand\red[1]{{\color{red}#1}}

\newcommand{\RQ}[1]{\textbf{RQ#1}}

\newcommand\owltopddl{Owl2PDDL\xspace}
\newcommand\acronym{PLANTA\xspace}
\newcommand{\CamaraApproach}{\textsc{Camara}\xspace}
\newcommand\plansys{PlanSys2\xspace}
\newcommand\tomasys{TOMASys\xspace}
\newcommand\pddltomasys{PDDL-TOMASys\xspace}

% suave uncertainties
\newcommand{\Uone}{\emph{U1}\xspace}
\newcommand{\Utwo}{\emph{U2}\xspace}
\newcommand{\Uthree}{\emph{U3}\xspace}

% navigation planta commands
\newcommand{\Sone}{\emph{S1}\xspace}
\newcommand{\Stwo}{\emph{S2}\xspace}
\newcommand{\Sthree}{\emph{S3}\xspace}
\newcommand{\Sfour}{\emph{S4}\xspace}
\newcommand{\Sfive}{\emph{S5}\xspace}


% define the colors to use. This is also the order how preferred they are
\definecolor{MidnightBlue}{HTML}{006795}
\definecolor{OliveGreen}{HTML}{3C8031}
\definecolor{BrickRed}{HTML}{B6321C}
\definecolor{Purple}{HTML}{800080}
\definecolor{Orange}{HTML}{F59527}
\newcommand{\greenColor}{OliveGreen} 	% encodes PLANTA
\newcommand{\blueColor}{MidnightBlue}	% encodes Camara et al
\newcommand{\redColor}{BrickRed}
\newcommand{\purpleColor}{Purple}
\newcommand{\orangeColor}{Orange}
% \newcommand{\whiteColor}{white}
% \newcommand\white[1]{{\color{white}#1}}

% tikz styles
\tikzstyle{lineDefault} = [
    line width=0.75pt
]

\tikzstyle{artifact} = [
lineDefault,
draw=black,
inner sep=3pt,
minimum width=1cm,
minimum height=0.5cm,
font=\sffamily,
fill=white!20
]

\tikzstyle{artifact_transparent} = [
lineDefault,
draw=black,
inner sep=3pt,
minimum width=1cm,
minimum height=0.5cm,
font=\sffamily,
]

\tikzstyle{process} = [
artifact,
rounded corners=0.2cm,
minimum height=0.3cm,
fill=black!5
]

\tikzstyle{tool} = [
    process,
    fill=\greenColor!20,
    draw=\greenColor
]

\tikzstyle{tool_3rd_party} = [
    process,
    fill=\blueColor!20,
    draw=\blueColor
]

\tikzstyle{system} = [
    process,
    fill=\purpleColor!5,
    draw=\purpleColor
]

\tikzstyle{runtime} = [
    artifact,
    fill=\blueColor!20,
    draw=\blueColor
]

\tikzstyle{arrow} = [
    line width=1pt,
    -Latex
]

\tikzstyle{arrow_dashed} = [
    line width=1pt,
    dashed,
    -Latex
]

\tikzstyle{bubble} = [
    ellipse,
    line width=1pt,
    draw
]

\tikzstyle{bubblePurple} = [
    bubble,
    fill=\purpleColor!20,
    draw=\purpleColor
]

\tikzstyle{bubbleGreen} = [
    bubble,
    fill=\greenColor!20,
    draw=\greenColor
]

\tikzstyle{bubbleBlue} = [
bubble,
fill=\blueColor!20,
draw=\blueColor
]

\newcommand{\gridNodeSize}{3pt}
\newcommand{\gridNodeLine}{0.75pt}	% should be grid size +2pt
\newcommand{\gridNodeSizePlus}{\gridNodeSize+\gridNodeLine+\gridNodeLine}

\tikzstyle{gridNode} = [
    minimum height=\gridNodeSize,
    minimum width=\gridNodeSize,
    inner sep=2.5pt,
    draw=black!70,
    line width=\gridNodeLine,
    fill=white,
    shape=circle
]

\tikzstyle{corridor} = [
    line width=\gridNodeLine,
    double,
    double distance=\gridNodeSize,
    draw=black!70
]


\tikzstyle{unsafe} = [
    decorate=true,
    decoration=snake,
    line cap=round
]

\tikzstyle{dark} = [
    draw=black!80!\blueColor!60!white
]

\tikzstyle{corridorDark} = [
    line width=\gridNodeSizePlus,
    dark
]

\tikzstyle{corridorUnsafe} = [
    unsafe,
    corridor
]

\tikzstyle{corridorDarkUnsafe} = [
    unsafe,
    line width=\gridNodeSizePlus,
    dark
]

\tikzstyle{goal} = [
    fill=\greenColor,
    text=white,
    shape=circle,
    draw=black!70,
    line width=\gridNodeLine
]
\tikzstyle{start} = [
    fill=\greenColor,
    text=white,
    shape=circle,
    draw=black!70,
    line width=\gridNodeLine
]

\tikzstyle{hideNode} = [
    fill=white,
    minimum height=\gridNodeSize+\gridNodeLine+7pt,
    minimum width=\gridNodeSize+\gridNodeLine+7pt,
    shape=circle
]

\pgfdeclarelayer{background}
\pgfsetlayers{background,main}



% Input a generated table file (e.g. tables/wilcoxon_analysis/*.tex) without its
% table float, caption, and label, so that it can be placed in a subtable or given
% a new caption. The generated files stay untouched for regeneration.
\newcommand{\inputtablebody}[1]{\begingroup
  \renewenvironment{table}[1][]{}{}%
  \renewcommand{\caption}[2][]{}%
  \renewcommand{\label}[1]{}%
  \input{#1}\endgroup} % chktex 27 -- argument is a file path

% pgf stuff for box-plots
\makeatletter
\pgfplotsset{
    boxplot prepared from table/.code={
        \def\tikz@plot@handler{\pgfplotsplothandlerboxplotprepared}%
        \pgfplotsset{
            /pgfplots/boxplot prepared from table/.cd,
            #1,
        }
    },
    /pgfplots/boxplot prepared from table/.cd,
    table/.code={\pgfplotstablecopy{#1}\to\boxplot@datatable},
    row/.initial=0,
    make style readable from table/.style={
        #1/.code={
            \pgfplotstablegetelem{\pgfkeysvalueof{/pgfplots/boxplot prepared from table/row}}{##1}\of\boxplot@datatable
            \pgfplotsset{boxplot/#1/.expand once={\pgfplotsretval}}
        }
    },
    make style readable from table=lower whisker,
    make style readable from table=upper whisker,
    make style readable from table=lower quartile,
    make style readable from table=upper quartile,
    make style readable from table=median,
    make style readable from table=lower notch,
    make style readable from table=upper notch
}
\makeatother